
\documentclass[10pt,a4paper]{article}
\usepackage[margin=0.6in]{geometry}
\usepackage[T2A]{fontenc}
\usepackage[utf8]{inputenc}
\usepackage[russian,english]{babel}
\usepackage{enumitem}
\usepackage[hidelinks]{hyperref}
\usepackage{microtype}

\setlength{\parindent}{0pt}
\setlength{\parskip}{0pt}
\setlength{\emergencystretch}{2em}
\setlist[itemize]{leftmargin=1.2em, itemsep=0pt, topsep=0pt}
\pagestyle{empty}

\newcommand{\cvsection}[1]{
  \vspace{3pt}
  {\large\textbf{#1}}\\[-4pt]
  \rule{\linewidth}{0.4pt}
}

\begin{document}
\raggedright

{\LARGE \textbf{У.Болорчулуун}}\\
Хэнтий, Монгол \;|\; +976 94972916 \;|\; \href{mailto:bolorchuluun.uuganbayar@gmail.com}{bolorchuluun.uuganbayar@gmail.com}\\
LinkedIn: \href{https://www.linkedin.com/in/bolroo-bolr0-179545355/}{linkedin.com/in/bolroo-bolr0-179545355} \;|\;
GitHub: \href{https://github.com/Blrch0n}{github.com/Blrch0n}

\cvsection{Мэргэжлийн зорилго}
Кибер аюулгүй байдлын чиглэлээр суралцаж буй оюутан; сүлжээний аюулгүй байдал, веб эмзэг байдлын үнэлгээ,
лог шинжилгээ, инцидентэд хариу арга хэмжээний суурь мэдлэгтэй. Практик лаборатори, веб систем хөгжүүлэлт,
багийн манлайллын туршлагадаа тулгуурлан дадлага эсвэл junior analyst түвшний боломж хайж байна.

\cvsection{Техникийн ур чадвар}
\begin{itemize}
  \item \textbf{Програмчлал:} Python, Java, C/C++, Bash
  \item \textbf{Аюулгүй байдал:} OWASP Top 10, Эмзэг байдлын үнэлгээ, Инцидентэд хариу арга хэмжээний суурь мэдлэг, SIEM суурь мэдлэг
  \item \textbf{Хэрэгслүүд:} Wireshark, Nmap, Burp Suite, Metasploit, Git, Linux
  \item \textbf{Сүлжээ:} TCP/IP, OSI загвар, Routing \& Switching-ийн суурь
  \item \textbf{Бусад:} Техникийн баримтжуулалт, Хувилбарын удирдлага, Багаар хамтран ажиллах
\end{itemize}

\cvsection{Техникийн төслүүд / Кибер лаборатори}
\textbf{Olympiad Management System} \hfill 2025.10 -- 2025.11\\
Вэб систем (олимпиадын үйл ажиллагааны менежмент)\\
GitHub: \href{https://github.com/Blrch0n/Olympiad-hugjuulelt}{github.com/Blrch0n/Olympiad-hugjuulelt}\\
Deployment: \href{https://olympiad.syscotech.club/}{olympiad.syscotech.club}\\
\begin{itemize}
  \item Олимпиадын бүртгэл, мэдээлэл удирдлага, хэрэглэгчийн веб хандалтад зориулсан онлайн систем хөгжүүлж, ашиглалтад бэлэн орчинд байршуулсан.
  \item Git/GitHub ашиглан хөгжүүлэлт, байршуулалт, системийн тасралтгүй ажиллагааг хангаж ажилласан.
  \item Ашигласан технологи: Next.js, TypeScript, Three.js, Firebase
\end{itemize}

\textbf{Network Traffic Analysis Lab (Wireshark)} \hfill 2026\\
\begin{itemize}
  \item Wireshark ашиглан сүлжээний урсгалыг барьж, TCP/UDP, DNS, HTTP протоколын хөдөлгөөнд шинжилгээ хийж, сэжигтэй хэв шинжийг ялган таньсан.
\end{itemize}

\textbf{Web Vulnerability Assessment Lab (Burp Suite \& OWASP Top 10)} \hfill 2026\\
\begin{itemize}
  \item Burp Suite ашиглан вебийн аюулгүй байдлын тест хийж, OWASP Top 10 ангиллаар эмзэг байдлын эрсдэлийг үнэлэн засварлах зөвлөмжийн жагсаалт бэлтгэсэн.
\end{itemize}

\cvsection{Ажлын туршлага / Манлайлал}
\textbf{Sys\&CoTech клуб, ШУТИС} -- Хөгжүүлэлтийн ахлагч \hfill 2025.09 -- Одоо\\
Клубын хуудас: \href{https://www.facebook.com/SysAndCoTech/}{facebook.com/SysAndCoTech}\\
\begin{itemize}
  \item 14 гишүүнтэй хөгжүүлэлтийн багийн ажлын төлөвлөлт, даалгавар хуваарилалт, хэрэгжилтийн хяналтыг удирдсан.
  \item Git ашигласан хувилбарын удирдлага, даалгаврын зохион байгуулалт, кодын хяналтын урсгалыг зохион байгуулсан.
\end{itemize}

\cvsection{Боловсрол}
\textbf{Шинжлэх Ухаан Технологийн Их Сургууль} \hfill 2024 -- 2027\\
\textit{Кибер аюулгүй байдал, Бакалавр}\\
Голч дүн: 3.14/4.00

\cvsection{Сертификат \& сургалт}
\begin{itemize}
  \item Cisco Networking Academy --- CCNA: Switching, Routing, and Wireless Essentials --- 2026 -- Одоо
  \item Cisco Networking Academy --- Junior Cybersecurity Analyst Career Path --- 2026 -- Одоо
  \item Pinecone Academy --- Developer Program --- 2023 -- 2024
\end{itemize}

\cvsection{Хэлний мэдлэг}
\begin{itemize}
  \item Монгол -- Төрөлх
  \item Англи -- Дунджаас дээш (техникийн уншлага, баримт бичиг боловсруулах)
\end{itemize}

\cvsection{Нэмэлт мэдээлэл}
\begin{itemize}
  \item Сонирхол: CTF, вэб эмзэг байдлын тест, реверс инженерчлэл, лог шинжилгээ
  \item Сайн дурын болон олон нийтийн үйл ажиллагаа: Цус хандивлах сайн дурын аянд оролцож, клубын төсөл, арга хэмжээнд идэвхтэй оролцдог
\end{itemize}

\end{document}
